Se desarrolló e implementó \textbf{UrbanTracker}, un sistema innovador de geolocalización en tiempo real orientado al transporte público urbano, que integra aplicaciones para usuarios, conductores y administradores en una plataforma unificada. El sistema cumple con los objetivos planteados de ofrecer información inmediata y precisa sobre la ubicación de autobuses a los pasajeros, y de facilitar la gestión operativa a los administradores mediante una vista global de la flota y herramientas administrativas digitales.

Entre las contribuciones principales de UrbanTracker, destaca su enfoque de \textbf{bajo costo y alta adaptabilidad}: al usar smartphones comunes como dispositivo de rastreo y tecnologías de software libre (WebSockets, MQTT, APIs web), se demuestra que incluso sistemas de transporte con recursos limitados pueden adoptar soluciones inteligentes de seguimiento. Asimismo, UrbanTracker aporta una \textbf{perspectiva integral} al abarcar las necesidades de distintos actores (usuario, conductor, administrador) en una sola plataforma, aspecto que lo diferencia de otras aplicaciones centradas solo en el pasajero.

En términos de desempeño, se logró una actualización de datos en tiempo real con frecuencias del orden de segundos y latencias muy reducidas, lo que excede las necesidades típicas del caso de uso y garantiza una experiencia fluida. Esto sienta las bases para mejorar la confianza de los usuarios en el transporte público, al reducir la incertidumbre sobre tiempos de llegada y permitir una mejor planificación de sus viajes diarios. Del lado administrativo, la posibilidad de monitorizar todos los vehículos activos en vivo puede traducirse en una operación más eficiente y reactiva, optimizando rutas y recursos en función de la demanda y las eventualidades (tráfico, averías, etc.).

Este trabajo también corrobora hallazgos de artículos de investigación sobre la eficacia de las arquitecturas pub/sub en contextos de IoT y rastreo vehicular, contribuyendo con una implementación específica que puede servir de referencia para otras ciudades o regiones interesadas en proyectos similares. UrbanTracker se presenta, en esencia, como un \textbf{caso de estudio exitoso} de transformación digital en el transporte público local.

Como trabajos futuros, se sugieren las siguientes acciones: (1) Realizar un despliegue piloto del sistema en un entorno real (por ejemplo, en una línea de buses urbana) para recopilar datos de uso real y retroalimentación de usuarios finales, ajustando la interfaz y funcionalidades de acuerdo a sus necesidades; (2) Incorporar algoritmos de estimación de tiempo de arribo de buses a paradas, utilizando la información histórica de recorridos y patrones de tráfico, lo que agregaría un valor significativo para los pasajeros; (3) Ampliar el soporte multiplataforma de la aplicación móvil (incluir iOS) y considerar una aplicación específica para pasajeros en móviles, complementando la web, para enviar notificaciones push de alertas o cambios en el servicio; (4) Profundizar en las medidas de seguridad y privacidad, asegurando cumplimiento con regulaciones de protección de datos y minimizando cualquier riesgo asociado al rastreo continuo (por ejemplo, anonimizar registros históricos y permitir a conductores o usuarios opciones de privacidad cuando aplique).

En conclusión, UrbanTracker demuestra cómo la integración de tecnologías de geolocalización, comunicación en tiempo real y desarrollo multiplataforma puede modernizar el servicio de transporte urbano, haciéndolo más \textbf{transparente, eficiente y centrado en el usuario}. Se espera que este proyecto sirva como base para futuras iniciativas de ciudades inteligentes en la región y que su implementación práctica genere beneficios tangibles en la calidad del transporte público y la satisfacción de sus usuarios.
